\documentclass[11pt,letterpaper,openany]{book}

% ──────────────────────────────────────────────
% Packages
% ──────────────────────────────────────────────
\usepackage[margin=1.2in]{geometry}
\usepackage{amsmath, amssymb, amsthm}
\usepackage{mathtools}
\usepackage{bm}
\usepackage{tikz}
\usepackage{pgfplots}
\pgfplotsset{compat=1.18}
\usetikzlibrary{arrows.meta, decorations.markings, calc, patterns, positioning, 3d}
\usepackage{tcolorbox}
\tcbuselibrary{theorems, breakable, skins}
\usepackage{hyperref}
\usepackage{enumitem}
\usepackage{fancyhdr}
\usepackage{titlesec}
\usepackage{epigraph}
\usepackage{xcolor}
\usepackage{booktabs}

% ──────────────────────────────────────────────
% Colors
% ──────────────────────────────────────────────
\definecolor{chapblue}{RGB}{0, 70, 130}
\definecolor{accentred}{RGB}{180, 40, 40}
\definecolor{softgray}{RGB}{245, 245, 245}
\definecolor{trajgreen}{RGB}{30, 130, 60}
\definecolor{darkgold}{RGB}{160, 120, 20}
\definecolor{purplehaze}{RGB}{100, 40, 140}

% ──────────────────────────────────────────────
% Theorem environments
% ──────────────────────────────────────────────
\newtcbtheorem[number within=chapter]{principle}{Mathematical Principle}{
  colback=chapblue!5, colframe=chapblue!80,
  fonttitle=\bfseries, separator sign={.~},
  breakable
}{princ}

\newtcbtheorem[number within=chapter]{laymansbox}{Intuition}{
  colback=trajgreen!5, colframe=trajgreen!80,
  fonttitle=\bfseries, separator sign={.~},
  breakable
}{intuition}

\newtcbtheorem[number within=chapter]{example}{Example}{
  colback=softgray, colframe=black!40,
  fonttitle=\bfseries, separator sign={.~},
  breakable
}{ex}

\newtcbtheorem[number within=chapter]{definition}{Definition}{
  colback=darkgold!5, colframe=darkgold!80,
  fonttitle=\bfseries, separator sign={.~},
  breakable
}{def}

\theoremstyle{plain}
\newtheorem{proposition}{Proposition}[chapter]
\newtheorem{theorem}{Theorem}[chapter]
\newtheorem{lemma}[theorem]{Lemma}

% ──────────────────────────────────────────────
% Chapter formatting
% ──────────────────────────────────────────────
\titleformat{\chapter}[display]
  {\normalfont\huge\bfseries\color{chapblue}}
  {\chaptertitlename\ \thechapter}{20pt}{\Huge}
\titlespacing*{\chapter}{0pt}{-20pt}{30pt}

% ──────────────────────────────────────────────
% Header/Footer
% ──────────────────────────────────────────────
\pagestyle{fancy}
\fancyhf{}
\fancyhead[LE,RO]{\thepage}
\fancyhead[RE]{\textit{Volume 0: The Mathematical Primer}}
\fancyhead[LO]{\textit{\nouppercase{\leftmark}}}
\renewcommand{\headrulewidth}{0.4pt}
\setlength{\headheight}{14.5pt}

% ──────────────────────────────────────────────
% Custom commands
% ──────────────────────────────────────────────
\newcommand{\state}{\bm{x}}
\newcommand{\control}{\bm{u}}
\newcommand{\traj}{\bm{\gamma}}
\newcommand{\manifold}{\mathcal{M}}
\newcommand{\configspace}{\mathcal{Q}}
\newcommand{\statespace}{\mathcal{X}}
\newcommand{\tube}{\mathcal{T}}
\newcommand{\Reals}{\mathbb{R}}
\newcommand{\dd}{\mathrm{d}}
\newcommand{\dt}{\,\dd t}
\newcommand{\norm}[1]{\left\lVert #1 \right\rVert}
\newcommand{\inner}[2]{\left\langle #1,\, #2 \right\rangle}
\newcommand{\se}{\mathfrak{se}(3)}
\newcommand{\so}{\mathfrak{so}(3)}
\newcommand{\SE}{\mathrm{SE}(3)}
\newcommand{\SO}{\mathrm{SO}(3)}
\newcommand{\screw}{\mathcal{S}}
\newcommand{\twist}{\mathcal{V}}
\newcommand{\wrench}{\mathcal{F}}

% ──────────────────────────────────────────────
% Title Page
% ──────────────────────────────────────────────
\title{%
  {\Huge\bfseries\color{chapblue} The Geometry of Motion}\\[0.4cm]
  {\Large Volume 0: The Mathematical Primer\\
  A Foundational Guide to Kinematics, Configuration, and State Space}\\[0.6cm]
  {\color{chapblue}\rule{\textwidth}{2pt}}
}
\author{{\Large D.\ Kraft}}
\date{}

\begin{document}

\maketitle
\thispagestyle{empty}
\cleardoublepage

\frontmatter
\chapter*{Preface}
\addcontentsline{toc}{chapter}{Preface}

Before we can begin to analyze the geometry of moving physical bodies or command them to navigate the world via control theory, we must establish a common language. The physics of movement happens in the real world, but our understanding of it happens entirely within mathematical spaces.

This primer, **Volume 0** of *The Geometry of Motion*, is designed to bridge the gap between basic calculus/algebra and the rigorous differential geometry required for modern nonlinear control theory. By intuitively introducing State Space, Configuration Manifolds, SO(3) Rotations, and Screw Theory, this text ensures that the "lay user" possesses the foundational armor needed to attack the subsequent volumes on Tangent Spaces (Volume I) and Trajectory Control (Volume II).

The philosophy of this primer maintains our overarching theme: mathematics should be driven by the physical reality of motion, not abstract algebra for its own sake.

\tableofcontents
\cleardoublepage

\mainmatter

% ==========================================================================
%  Chapter 1: The Transition to State Space
% ==========================================================================
\chapter{The Transition to State Space}
\label{ch:state_space}

\epigraph{%
  You cannot control a machine by staring at the machine. You control it by staring at the mathematical space that completely defines it.
}{}

\section{Introduction: The Physics Abstraction}

When a layman pictures a robotic arm throwing a ball or a drone correcting its pitch in a crosswind, they visualize physical bodies acting in 3-dimensional Euclidean space (the $X, Y, Z$ coordinates of our physical reality). 

But control theorists and kinematic engineers almost never look at the $X, Y, Z$ space. They work in a wholly abstract, higher-dimensional mathematical realm known as \textbf{State Space}. State space is arguably the most fundamental translation required in control theory. If you do not understand state space, nothing else in the subsequent volumes of *The Geometry of Motion* will make sense.

\begin{laymansbox}{What is State Space?}{statespace_intuition}
  Imagine taking a photograph of a physical system. A photograph captures position, but not motion (you cannot tell if a car in a photo is parked or moving 100 mph). State Space is the mathematical equivalent of the ultimate, complete "photograph." A single point in State Space contains absolutely everything you need to know about the system to perfectly predict its future (assuming you know the laws of physics dictating it).
\end{laymansbox}

\section{Constructing a State Vector}

Let us build a state space from the ground up using the simplest dynamical system possible: a 1-Dimensional point mass (like a train on a straight track).

We want to predict what the train will do. To do this, we need to know its position, $x$. Is $x$ the "state" of the system?

No. If we only know that the train is at $x = 100$ meters, we cannot predict where it will be 1 second from now. It could be moving right, moving left, or stationary. 

Therefore, to complete our "photograph," we must also know its velocity, $v$ (often denoted as the derivative of position, $\dot{x}$). If we know $(x, \dot{x})$, we know its exact location and its exact speed.

Do we need to know acceleration ($\ddot{x}$) as part of the state? According to Newton's Second Law: $F = m\ddot{x}$. Acceleration is not an independent property of the train; it is completely driven by the external forces applied at that instant. Therefore, if we know the forces (the Control Inputs, $\control$), we can compute the acceleration. The state itself is fully defined by just position and velocity.

\begin{definition}{The State Vector $\state$}{statevector}
  The state vector $\state$ is the minimum set of mathematical variables completely describing the condition of a dynamic system at a specific point in time. 
  
  For our train, it is a 2-dimensional vector:
  \begin{equation}
      \state = \begin{bmatrix} x_1 \\ x_2 \end{bmatrix} = \begin{bmatrix} \text{position} \\ \text{velocity} \end{bmatrix} = \begin{bmatrix} q \\ \dot{q} \end{bmatrix}
  \end{equation}
\end{definition}

\section{State Space Representations}

Because the state vector $\state$ contains $n$ variables, it naturally lives in an $n$-dimensional geometry: the State Space, denoted as $\statespace \in \Reals^n$.

An astonishing realization for beginners is that when a dynamic system (like an airplane or a robotic arm) moves through physical space over time, its entire sprawling motion is represented as a single, unified \textbf{1-dimensional curve} snaking through $n$-dimensional state space. This curve is called the system's \textbf{Trajectory} $\traj$.

\subsection{Standard Form Differential Equations}

The immense power of state space is that it allows us to convert complex, higher-order physical differential equations into simple, first-order vector math.

Consider a simple mass-spring-damper system governed by a second-order differential equation:
\begin{equation}
  m\ddot{q} + c\dot{q} + kq = u
\end{equation}
where $q$ is position, $m$ is mass, $c$ is damping, $k$ is stiffness, and $u$ is our control force.

By defining our state vector as $\state = [x_1, x_2]^T = [q, \dot{q}]^T$, we can rewrite this physical equation as purely mathematical first-order derivatives:

\begin{align}
  \dot{x}_1 &= x_2 \\
  \dot{x}_2 &= \frac{1}{m}(u - cx_2 - kx_1)
\end{align}

This is called the \textbf{State Space Representation}. It is universally written in nonlinear control literature as:

\begin{principle}{The Fundamental Equation of Control Theory}{fund_eq}
  Every autonomous control system can be written as a first-order vector differential equation mapping current state and current control commands to the velocity of the state vector:
  \begin{equation}
      \bm{\dot{x}} = \bm{f}(\state, \control, t)
  \end{equation}
  This vector field $\bm{f}$ dictates how the system "flows" through the abstract geometry of the state space. 
\end{principle}

Whether you are controlling a single gear or a 15-DOF musculoskeletal golf swing (as discussed in Volume II), you are always, fundamentally, trying to shape the trajectory $\state(t)$ by manipulating the vector field $\bm{f}$ using your control inputs $\control$.

\section{Summary}

State space fundamentally decouples physics from geometry. We no longer treat a system as physical objects bumping into each other in 3D space. We treat it as a single point—a cursor—moving through hyper-dimensional coordinates based on a vector field. 

But what defines those coordinates? How do we mathematically decide what the angles and positions of highly complex mechanisms are? For that, we turn to Configuration Space.

\include{chapters/ch02_configuration}
% ==========================================================================
%  Chapter 3: Rotations, Quaternions, and SO(3)
% ==========================================================================
\chapter{Rotations, Quaternions, and SO(3)}
\label{ch:rotations}

\epigraph{%
  The tragedy of Euclidean space is that it works perfectly until you spin.
}{}

\section{The Coordinate Frame}

When we model movement or control robotic systems, we are not looking at the physical object itself; we are tracking a \textbf{Coordinate Frame} rigidly attached to the object, often denoted as the Body Frame ($B$). 

To know where the Body Frame $B$ is relative to the absolute stationary universe—the Spatial (or World) Frame ($S$)—we must define two mathematical properties:
1. The **translation** vector $\bm{p} \in \Reals^3$ pointing from the origin of $S$ to the origin of $B$. 
2. The **rotation** of the axes of $B$ relative to the axes of $S$.

Unlike the straightforward $(X, Y, Z)$ translation vector, rotation is profoundly mathematically problematic.

\section{Euler Angles and Gimbal Lock}

The most intuitive way to describe the orientation of an aircraft, drone, or robotic arm is to split rotation into three successive angles around axes, usually denoted as Roll, Pitch, and Yaw ($X, Y, Z$). 

This is known as an \textbf{Euler Angle} representation: $\bm{\theta} = [\theta_1, \theta_2, \theta_3]^T$. While highly intuitive for a human pilot, Euler Angles harbor a fatal mathematical flaw known as \textbf{Gimbal Lock}.

\begin{laymansbox}{Gimbal Lock}{gimballock}
  If an airplane pitches perfectly $90^\circ$ straight up into the sky, its Roll axis (pointing out its nose) now perfectly aligns with its original Yaw axis (pointing perfectly down towards the Earth). 
  
  Mathematically, two independent degrees of freedom have become linearly dependent—they merged into the identical physical rotation axis. A degree of freedom is literally lost during the calculation. If you attempt to control the plane passing through this "Singularity", your matrices become non-invertible and your computer instantly crashes mapping a physical impossibility.
\end{laymansbox}

Since our control theory and geometric motion algorithms fundamentally rely on continuous trajectories twisting through all possible spatial angles (like chaotic golf club arcs or throwing motions), we must completely abandon Euler angles. We cannot afford mathematical singularities.

\section{The Special Orthogonal Group: SO(3)}

To map 3D orientation without singularities in all continuous space, we define a rotation by looking at the purely physical unit vectors defining the Body Frame $B$: the $\hat{x}_b, \hat{y}_b, \hat{z}_b$ axes.

If we write these three body axes explicitly in the coordinates of the Space Frame $S$, we can stack them vertically into a $3 \times 3$ matrix:
\begin{equation}
   R = \begin{bmatrix} | & | & | \\ \hat{x}_b & \hat{y}_b & \hat{z}_b \\ | & | & | \end{bmatrix} \in \Reals^{3 \times 3}
\end{equation}

This is a \textbf{Rotation Matrix}. 

Because $\hat{x}_b, \hat{y}_b,$ and $\hat{z}_b$ are all length $1$ (unit magnitude) and perfectly orthogonal (all $90^\circ$ apart), the columns are an orthonormal basis. The strict mathematical definitions that guarantee this orthonormality are:
1. $R^T R = I$ (The transpose equals the inverse)
2. $\det(R) = 1$ (The determinant is 1 to preserve right-handed orientation—no reflecting through a mirror)

\begin{principle}{The SO(3) Manifold}{so3}
  The continuous, infinite geometric space of all possible valid $3 \times 3$ rotation matrices satisfying $R^T R = I$ and $\det(R) = 1$ is a mathematical \textbf{manifold} called the \textbf{Special Orthogonal Group in 3 Dimensions}, denoted as $\SO$. 

  If a Configuration Space $\configspace$ is simply a rigid body rotating around a point, we say $\configspace = \SO$.
\end{principle}

\section{Unit Quaternions $\so$}

While $\SO$ completely prevents Gimbal Lock and provides pure, continuous rotation matrices, tracking 9 independent variables inside a $3 \times 3$ grid is wildly computationally inefficient for an engine updating a 15-DOF human model at 1,000 Hertz. 

To solve this, William Rowan Hamilton invented \textbf{Quaternions} ($\mathcal{H}$) in 1843. Rather than tracking a full Cartesian body frame $R \in \SO$, any 3D rotation can be flawlessly described as a combination of precisely two things:
1. A single unit vector $\hat{\bm{\omega}}$ defining the solitary physical axis of rotation.
2. An angle $\theta$ denoting how far to spin around that axis.

A Unit Quaternion $q$ mathematically encodes these into 4 variables representing an exponential scalar and vector map:
\begin{equation}
  q = \begin{bmatrix} q_0 \\ q_1 \\ q_2 \\ q_3 \end{bmatrix} = \begin{bmatrix} \cos(\theta/2) \\ \hat{\bm{\omega}}_x \sin(\theta/2) \\ \hat{\bm{\omega}}_y \sin(\theta/2) \\ \hat{\bm{\omega}}_z \sin(\theta/2) \end{bmatrix} \in \Reals^4
\end{equation}

By simply enforcing that the sum of the squares of these 4 variables equals 1 ($q_0^2 + q_1^2 + q_2^2 + q_3^2 = 1$), a Unit Quaternion provides a singularity-free, ultra-fast algebraic mapping of $\SO$.

Throughout *The Geometry of Motion*, specifically when interacting with the codebase bridges connecting algorithms like `Pinocchio` or `UpstreamDrift` (which natively uses MuJoCo physics engines), all 3D rotational mechanics explicitly bypass Euler Angles and map directly into Quaternions and $\SO$ Rotation Matrices.

To merge Translation and Rotation into a single unified control framework for physical bodies mapping the full world, we must upgrade from $3 \times 3$ Rotation Matrices to the magnificent $4 \times 4$ space of $\SE$.

\include{chapters/ch04_screw_axes}

\backmatter
\bibliographystyle{plain}
\bibliography{../geometry_of_motion}

\end{document}
